\hypertarget{regular-expressions-engine-internals-re-sre_compile}{%
\chapter{\texorpdfstring{Regular Expressions Engine Internals
(\texttt{re},
\texttt{sre\_compile})}{Regular Expressions Engine Internals (re, sre\_compile)}}\label{regular-expressions-engine-internals-re-sre_compile}}

Regular expressions are a language within a language. While most
developers use them as black boxes, the CPython
\passthrough{\lstinline!re!} module is a sophisticated engine that
translates patterns into a custom bytecode executed by a specialized
virtual machine.

\hypertarget{the-sre-engine-architecture}{%
\subsection{\texorpdfstring{42.1 The \texttt{sre} Engine
Architecture}{42.1 The sre Engine Architecture}}\label{the-sre-engine-architecture}}

CPython's regex engine is called \textbf{sre} (Secret Rabbit Engine). It
is a \textbf{backtracking-based NFA} (Non-deterministic Finite
Automaton) engine.

\hypertarget{compilation-to-sre-bytecode}{%
\subsubsection{\texorpdfstring{1. Compilation to \texttt{sre}
Bytecode}{1. Compilation to sre Bytecode}}\label{compilation-to-sre-bytecode}}

When you call \passthrough{\lstinline!re.compile(pattern)!}, the
following happens: 1. \textbf{Parsing}: The \passthrough{\lstinline!re!}
module (in Python) parses the pattern string into a tree of tokens. 2.
\textbf{Optimization}: It performs optimizations like merging adjacent
literal characters into single ``string'' match commands. 3.
\textbf{Code Generation}: The \passthrough{\lstinline!sre\_compile!}
module translates this tree into a sequence of integer-based opcodes.

You can actually see these opcodes using the undocumented
\passthrough{\lstinline!re.purge()!} and looking at the compiled
object's \passthrough{\lstinline!.code!} attribute, or by setting
\passthrough{\lstinline!re.DEBUG!} during compilation:

\begin{lstlisting}[language=Python]
import re
re.compile("a(b|c)d", re.DEBUG)
\end{lstlisting}

\emph{Output (Simplified):}

\begin{lstlisting}
literal 97
subpattern 1
    branch
        literal 98
    or
        literal 99
literal 100
\end{lstlisting}

\hypertarget{the-matcher-vm-sre_lib.h}{%
\subsubsection{\texorpdfstring{2. The Matcher VM
(\texttt{sre\_lib.h})}{2. The Matcher VM (sre\_lib.h)}}\label{the-matcher-vm-sre_lib.h}}

The actual matching happens in C
(\passthrough{\lstinline!Modules/\_sre/sre.c!}). It is a recursive
function that takes the bytecode and the input string. *
\textbf{Backtracking}: When a branch fails (e.g., in
\passthrough{\lstinline!(b|c)!}), the engine ``backtracks'' to the last
save point and tries the next alternative. * \textbf{Performance Trap}:
Because it uses backtracking, ``catastrophic backtracking'' (exponential
time complexity) is possible with certain nested quantifiers.

\hypertarget{modern-enhancements-and-python-3.11}{%
\subsection{42.2 Modern Enhancements and Python
3.11+}\label{modern-enhancements-and-python-3.11}}

In Python 3.11, the \passthrough{\lstinline!re!} engine received
significant performance boosts. The ``atomic grouping''
(\passthrough{\lstinline!(?>...)!}) and possessive quantifiers
(\passthrough{\lstinline!*+!}, \passthrough{\lstinline!++!}) were added,
allowing developers to explicitly disable backtracking for specific
subpatterns, protecting against ReDoS (Regular Expression Denial of
Service) attacks.

\begin{center}\rule{0.5\linewidth}{0.5pt}\end{center}
